\section*{अध्याय \ref{ch:g4:animal-reproduction} --- जंतु कैसे जनन करते हैं}

\begin{solution}{exo:g4:animal-reproduction:1}
एक \omterm{def:g4:animal-reproduction:sexual}{शुक्राणु} का एक \omterm{def:g4:animal-reproduction:sexual}{अंडाणु} से जुड़ना, जिससे नया नन्हा \omterm{def:g1:animals-around-us:animal}{जंतु} शुरू होता है।
\omterm{def:g4:animal-reproduction:sexual}{अंडाणु} \omterm{def:g4:animal-reproduction:sexual}{मादा} से आता है, और \omterm{def:g4:animal-reproduction:sexual}{शुक्राणु} \omterm{def:g4:animal-reproduction:sexual}{नर} से।
\end{solution}

\begin{solution}{exo:g4:animal-reproduction:2}
\omterm{def:g4:animal-reproduction:ovivivi}{अंडज}: \omterm{def:g4:animal-reproduction:sexual}{मादा} \omterm{def:g2:animals-grow-up:born}{अंडे} देती है और बच्चा \omterm{def:g2:animals-grow-up:born}{अंडे} के भीतर, उसके शरीर के बाहर बढ़कर
पूरा होता है --- मुर्ग़ी, मेंढक (और ट्राउट, छिपकलियाँ, घोंघे, \omterm{prop:g3:sorting-living-things:groups}{कीट} भी)।
\omterm{def:g4:animal-reproduction:ovivivi}{जरायुज}: बच्चा माँ के भीतर बढ़ता है, उसी से पोषण पाता है, और \omterm{def:g2:animals-grow-up:born}{जीवित जन्म}
लेता है --- घोड़ी, ह्वेल (लगभग सारे \omterm{prop:g3:sorting-living-things:groups}{स्तनधारी})।
\end{solution}

\begin{solution}{exo:g4:animal-reproduction:3}
\omterm{def:g4:animal-reproduction:sexual}{नर} \omterm{def:g4:animal-reproduction:sexual}{मादाओं} को तालाब बुलाने के लिए गा रहे थे; फिर हर जोड़ी ने अपनी
कोशिकाएँ पानी में छोड़ीं, जहाँ \omterm{def:g4:animal-reproduction:sexual}{शुक्राणु} \omterm{def:g2:animals-grow-up:born}{अंडों} से मिले --- यानी \omterm{def:g4:animal-reproduction:sexual}{निषेचन};
और लुआब में लिपटे निषेचित \omterm{def:g2:animals-grow-up:born}{अंडे} बढ़कर जून के \omterm{ex:g2:animals-grow-up:frog}{टैडपोल} बन गए।
\end{solution}

\begin{solution}{exo:g4:animal-reproduction:4}
ट्राउट का \omterm{def:g4:animal-reproduction:sexual}{निषेचन} पानी में, \omterm{def:g4:animal-reproduction:sexual}{मादा} के बाहर होता है --- \omterm{def:g2:animals-grow-up:born}{अंडे} और \omterm{def:g4:animal-reproduction:sexual}{शुक्राणु}
साथ छोड़े जाते हैं। मुर्ग़ी का उसके शरीर के भीतर, \omterm{prop:g4:animal-reproduction:where}{समागम} के दौरान होता
है। ज़मीन पर खुली कोशिकाएँ सूख जातीं; पानी में वे बाहर भी सुरक्षित मिल
सकती हैं।
\end{solution}

\begin{solution}{exo:g4:animal-reproduction:5}
चूज़ा मुर्ग़ी के बाहर खोलदार \omterm{def:g2:animals-grow-up:born}{अंडे} में बढ़ता है और उसी के साथ भरी ज़र्दी
पर जीता है। बछेड़ा ग्यारह महीने घोड़ी के भीतर बढ़ता है, लगातार उसके
शरीर से पोषण पाता है, और \omterm{def:g2:animals-grow-up:born}{जीवित जन्म} लेता है।
\end{solution}

\begin{solution}{exo:g4:animal-reproduction:6}
प्रजाति जितने ज़्यादा बच्चे पैदा करती है, हर एक को उतनी ही कम देखभाल
मिलती है; जितने कम, उतनी ही ज़्यादा। सिरे: कॉड मछली के लाखों बिना
देखभाल वाले \omterm{def:g2:animals-grow-up:born}{अंडे}, और हथिनी का एक दशक तक रखवाला गया अकेला बच्चा।
\end{solution}

\begin{solution}{exo:g4:animal-reproduction:7}
क्योंकि सफलता के लिए बस इतना ज़रूरी है कि औसतन इतने बच्चे बच जाएँ कि
जनकों की जगह ले सकें। लाखों \omterm{def:g2:animals-grow-up:born}{अंडों} में से कुछ न कुछ हमेशा समुद्र के
ख़तरों से निकल ही जाते हैं --- वह विशाल संख्या \emph{ही} रक्षा है।
\end{solution}

\begin{solution}{exo:g4:animal-reproduction:8}
कम बच्चे (हज़ारों के मुक़ाबले पाँच कम ही हैं), माँ के भीतर बढ़े, \omterm{def:g2:animals-grow-up:born}{दूध}
पिलाए और रखवाली किए गए: बिल्ली साफ़ तौर पर कम-और-देखभाल-वाले पक्ष में
है --- आम तौर पर \omterm{prop:g3:sorting-living-things:groups}{स्तनधारियों} की तरह।
\end{solution}

\begin{solution}{exo:g4:animal-reproduction:9}
कि उसे \emph{हर} जनक से कुछ न कुछ मिलता है --- बच्चा अकेली अपनी माँ की
नक़ल नहीं है। बछेड़ा जो बनता है, उसमें दोनों जनकों का योगदान है।
\end{solution}

\begin{solution}{exo:g4:animal-reproduction:10}
छोटा गुच्छा --- कुछ दर्जन या कुछ सौ \omterm{def:g2:animals-grow-up:born}{अंडे}, लाखों नहीं। देखभाल और संख्या
का सौदा चलता है: जो पिता पहरा देता है, पानी हिलाता है और रक्षा करता
है, वह छोटे गुच्छे को उन ख़तरों से पार निकाल सकता है जिनके लिए बिना
पहरे के लाखों \omterm{def:g2:animals-grow-up:born}{अंडे} चाहिए होते।
\end{solution}

\begin{solution}{exo:g4:animal-reproduction:11}
\omterm{def:g4:animal-reproduction:ovivivi}{अंडज} --- और बिना देखभाल के: गाड़ा हुआ गुच्छा रेत की गरमाहट के भरोसे
छोड़ दिया जाता है, और निकले हुए बच्चे अकेले ही समुद्र की ओर दौड़ते हैं
(उनकी बड़ी संख्या ही उनकी रणनीति है)। उसे ज़मीन पर ही \omterm{def:g2:animals-grow-up:born}{अंडे} देने पड़ते
हैं क्योंकि \omterm{prop:g4:classifying-animals:classes}{सरीसृप} के \omterm{def:g2:animals-grow-up:born}{अंडे} खोलदार होते हैं और हवा में साँस लेते हैं
--- समुद्र में दिए जाने पर वे डूब जाते। नियम उसकी आदतें नहीं, उसके
वर्ग की ख़ासियतें तय करती हैं।
\end{solution}
